\centerline{\bf \Huge Метод математической индукции II}
\title{Метод математической индукции II}
\begin{flushright}
    \scriptsize{...Это индукция. То, что ты видишь, реально.\\ 
    Тебе ни за что не достичь того момента, когда твои действия станут реальны.\\
    Такова сила математической индукции.}
\end{flushright}

\textbf{Основная схема математической индукции на примере задачи:}

\par \textit{Пусть доминошки выставлены на ребро друг за другом. Нам необходимо доказать следующий факт: "Если толкнуть 1-ю доминошку, то для любого сколь угодно большого N доминошка с номером N когда-нибудь упадёт".}

\par\textbf{Для доказательства понадобится установить верность 2-х утверждений:}
\textit{1) База: мы можем толкнуть первую доминошку и она упадёт.}\\
\textit{2) Переход: Если подает доминошка с номером k, то она толкает доминошку с номером (k + 1)}\\

1. Докажите неравенство: $2^{n}>n$ для любых натуральных n.\\

2. Добрый Очир купил Лососям бесконечную пиццу (имеется в виду бесконечная плоскость). Так как у Лососей плавники и ровно разрезать пиццу не получается, они просто провели n хаотичных прямых разрезов. После запекания теста настала пора раскладывать ингридиенты $-$ колбакси и помидорки. Докажите, что Лососи могут разложить ингридиенты так, что кусок пиццы с колбаской граничит по прямой с куском пиццы с помидоркой и наоборот.\\

3. Докажите тождество: $(1-\dfrac{1}{4})(1-\dfrac{1}{9})\dots(1-\dfrac{1}{n^{2}})=\dfrac{n+1}{2n}$\\

$4^{*}$. Обозначим через $F_1, F_2, F_3, F_4, \dots$ последовательность чисел Фибоначчи. (По определению $F_1 = F_2 = 1$, а $F_k = F_{k-1} + F_{k-2}$ для любого $k>3$.) Докажите, что \\
\begin{center}
  $F_{mn}=F_mF_{n+1}+F_{m-1}F_n$
\end{center}
\\
для любых $m, n \geq 2$